\poczatekwykladu{3}
\section{Ideały i filtry zbiorów}
\label{sec:idealy-filtry}

W tej części wszystkie rozważane zbiory są podzbiorami ustalonego zbioru
$X$, a dopełnienia bierzemy względem $X$.

\begin{definicja}[Ideał zbiorów]
\label{def:ideal}
Rodzinę $\mathcal{I}\subseteq\mathcal{P}(X)$ nazywamy ideałem zbiorów,
jeżeli spełnia następujące warunki:
\begin{enumerate}
  \item $\varnothing\in\mathcal{I}$ oraz $X\notin\mathcal{I}$;
  \item jeżeli $A,B\in\mathcal{I}$, to $A\cup B\in\mathcal{I}$;
  \item jeżeli $A\subseteq B$ i $B\in\mathcal{I}$, to $A\in\mathcal{I}$.
\end{enumerate}
Ostatni warunek oznacza, że każdy podzbiór elementu ideału również
należy do tego ideału.
\end{definicja}

\begin{przyklad}[Ideał zbiorów skończonych]
\label{prz:ideal-skonczone}
Dla $X=\R$ rodzina
\begin{equation}
  \mathcal{I}_{\mathrm{fin}}
  =\{A\subseteq\R:A\text{ jest skończony}\}
  \label{eq:ideal-skonczone}
\end{equation}
jest ideałem zbiorów.
\end{przyklad}

\begin{proof}
Sprawdzamy kolejno warunki definicji~\ref{def:ideal}.
\begin{enumerate}
  \item Zbiór pusty jest skończony, więc
  $\varnothing\in\mathcal{I}_{\mathrm{fin}}$.
  Zbiór $\R$ jest nieskończony, zatem $\R\notin\mathcal{I}_{\mathrm{fin}}$.
  \item Jeżeli $A,B\in\mathcal{I}_{\mathrm{fin}}$, to oba zbiory są
  skończone. Ich suma jest skończona, gdyż
  $|A\cup B|\leq|A|+|B|<\infty$.
  Stąd $A\cup B\in\mathcal{I}_{\mathrm{fin}}$.
  \item Jeżeli $A\subseteq B$ i $B\in\mathcal{I}_{\mathrm{fin}}$,
  to $|A|\leq|B|<\infty$. Zatem $A$ jest skończony i należy
  do $\mathcal{I}_{\mathrm{fin}}$.
\end{enumerate}
\end{proof}

\begin{definicja}[Filtr zbiorów]
\label{def:filtr}
Rodzinę $\mathcal{F}\subseteq\mathcal{P}(X)$ nazywamy filtrem zbiorów,
jeżeli spełnia następujące warunki:
\begin{enumerate}
  \item $X\in\mathcal{F}$ oraz $\varnothing\notin\mathcal{F}$;
  \item jeżeli $A,B\in\mathcal{F}$, to $A\cap B\in\mathcal{F}$;
  \item jeżeli $A\subseteq B\subseteq X$ i $A\in\mathcal{F}$,
  to $B\in\mathcal{F}$.
\end{enumerate}
Ostatni warunek oznacza, że każdy nadzbiór elementu filtru, zawarty
w $X$, również należy do tego filtru.
\end{definicja}

\begin{przyklad}[Filtr zbiorów o skończonym dopełnieniu]
\label{prz:filtr-dopelnienia-skonczone}
Dla $X=\R$ rodzina
\begin{equation}
  \mathcal{F}_{\mathrm{cofin}}
  =\{A\subseteq\R:A^c\text{ jest skończony}\}
  \label{eq:filtr-dopelnienia-skonczone}
\end{equation}
jest filtrem zbiorów.
\end{przyklad}

\begin{proof}
Sprawdzamy kolejno warunki definicji~\ref{def:filtr}.
\begin{enumerate}
  \item Mamy $\R^c=\varnothing$, więc
  $\R\in\mathcal{F}_{\mathrm{cofin}}$.
  Natomiast $\varnothing^c=\R$ jest nieskończony, zatem
  $\varnothing\notin\mathcal{F}_{\mathrm{cofin}}$.
  \item Jeżeli $A,B\in\mathcal{F}_{\mathrm{cofin}}$, to
  $A^c$ i $B^c$ są skończone. Z prawa de Morgana otrzymujemy
  \[
    (A\cap B)^c=A^c\cup B^c.
  \]
  Prawa strona jest sumą dwóch zbiorów skończonych, więc jest skończona.
  Stąd $A\cap B\in\mathcal{F}_{\mathrm{cofin}}$.
  \item Jeżeli $A\subseteq B\subseteq\R$ i
  $A\in\mathcal{F}_{\mathrm{cofin}}$, to $B^c\subseteq A^c$.
  Ponieważ $A^c$ jest skończony, również $B^c$ jest skończony.
  Zatem $B\in\mathcal{F}_{\mathrm{cofin}}$.
\end{enumerate}
\end{proof}

\begin{fakt}[Dualność ideałów i filtrów]
\label{fakt:dualnosc-ideal-filtr}
Rodzina $\mathcal{I}\subseteq\mathcal{P}(X)$ jest ideałem wtedy
i tylko wtedy, gdy rodzina
\begin{equation}
  \mathcal{I}^{\ast}=\{A^c:A\in\mathcal{I}\}
  \label{eq:rodzina-dualna}
\end{equation}
jest filtrem.
\end{fakt}

\begin{proof}
Niech $\mathcal{I}$ będzie ideałem. Z warunków
$\varnothing\in\mathcal{I}$ i $X\notin\mathcal{I}$ otrzymujemy
$X\in\mathcal{I}^{\ast}$ i $\varnothing\notin\mathcal{I}^{\ast}$.
Jeżeli $C,D\in\mathcal{I}^{\ast}$, to $C^c,D^c\in\mathcal{I}$, więc
\[
  (C\cap D)^c=C^c\cup D^c\in\mathcal{I}.
\]
Zatem $C\cap D\in\mathcal{I}^{\ast}$.
Jeżeli $C\subseteq D\subseteq X$ i $C\in\mathcal{I}^{\ast}$,
to $D^c\subseteq C^c\in\mathcal{I}$. Zamknięcie ideału na branie
podzbiorów daje $D^c\in\mathcal{I}$, czyli $D\in\mathcal{I}^{\ast}$.

Odwrotnie, niech $\mathcal{I}^{\ast}$ będzie filtrem.
Mamy $\varnothing\in\mathcal{I}$ i $X\notin\mathcal{I}$.
Jeżeli $A,B\in\mathcal{I}$, to
$(A\cup B)^c=A^c\cap B^c\in\mathcal{I}^{\ast}$,
więc $A\cup B\in\mathcal{I}$.
Jeżeli $A\subseteq B$ i $B\in\mathcal{I}$, to
$B^c\subseteq A^c$ i $B^c\in\mathcal{I}^{\ast}$.
Z warunku dotyczącego nadzbiorów w filtrze wynika
$A^c\in\mathcal{I}^{\ast}$, a więc $A\in\mathcal{I}$.
\end{proof}

\section{Pi-układy i lambda-układy}
\label{sec:pi-lambda-uklady}

\begin{definicja}[Pi-układ]
\label{def:pi-uklad}
Rodzinę $\mathcal{P}\subseteq\mathcal{P}(X)$ nazywamy pi-układem
($\pi$-układem), jeżeli jest niepusta i zamknięta na przekroje dwóch
zbiorów, to znaczy
\[
  A,B\in\mathcal{P}\quad\Longrightarrow\quad A\cap B\in\mathcal{P}.
\]
\end{definicja}

\begin{przyklad}[Rodziny będące pi-układami]
W przestrzeni $X=\R$ następujące rodziny są pi-układami:
\begin{enumerate}
  \item rodzina $\mathcal{I}_{\mathrm{fin}}$ zbiorów skończonych;
  \item rodzina $\mathcal{F}_{\mathrm{cofin}}$ zbiorów o skończonym dopełnieniu;
  \item rodzina $\mathcal{J}$ wszystkich przedziałów w $\R$ wraz
  ze zbiorem pustym; dopuszczamy także przedziały nieograniczone.
\end{enumerate}
\end{przyklad}

\begin{proof}
Każda z tych rodzin jest niepusta. Przekrój dwóch zbiorów skończonych
jest skończony, gdyż jest podzbiorem każdego z nich.
Zamknięcie rodziny $\mathcal{F}_{\mathrm{cofin}}$ na przekroje
sprawdziliśmy w przykładzie~\ref{prz:filtr-dopelnienia-skonczone}.
Przekrój dwóch przedziałów jest przedziałem albo zbiorem pustym,
więc również należy do $\mathcal{J}$.
\end{proof}

\begin{definicja}[Lambda-układ]
\label{def:lambda-uklad}
Rodzinę $\mathcal{L}\subseteq\mathcal{P}(X)$ nazywamy lambda-układem
($\lambda$-układem), jeżeli:
\begin{enumerate}
  \item $X\in\mathcal{L}$;
  \item jeżeli $A,B\in\mathcal{L}$ i $A\subseteq B$,
  to $B\setminus A\in\mathcal{L}$;
  \item jeżeli $A_n\in\mathcal{L}$ dla każdego $n\geq1$ oraz
  $A_1\subseteq A_2\subseteq\cdots$, to
  $\bigcup_{n=1}^{\infty}A_n\in\mathcal{L}$.
\end{enumerate}
\end{definicja}

\begin{twierdzenie}[Dynkina--Sierpińskiego]
\label{tw:pi-lambda-sigma}
Rodzina zbiorów, która jest jednocześnie pi-układem i lambda-układem,
jest sigma-ciałem.
\end{twierdzenie}

\begin{twierdzenie}[O pi-układzie i lambda-układzie]
\label{tw:pi-lambda}
Jeżeli $\mathcal{P}$ jest pi-układem, $\mathcal{L}$ jest lambda-układem
oraz $\mathcal{P}\subseteq\mathcal{L}$, to
\begin{equation}
  \sigma(\mathcal{P})\subseteq\mathcal{L}.
  \label{eq:pi-lambda}
\end{equation}
\end{twierdzenie}

\section{Rodziny monotoniczne}
\label{sec:rodziny-monotoniczne}

\begin{definicja}[Rodzina monotoniczna]
\label{def:rodzina-monotoniczna}
Rodzinę $\mathcal{M}\subseteq\mathcal{P}(X)$ nazywamy rodziną
monotoniczną, jeżeli:
\begin{enumerate}
  \item jest zamknięta na przeliczalne sumy ciągów wstępujących:
  jeżeli $A_n\in\mathcal{M}$ i $A_1\subseteq A_2\subseteq\cdots$, to
  \[
    \bigcup_{n=1}^{\infty}A_n\in\mathcal{M};
  \]
  \item jest zamknięta na przeliczalne przekroje ciągów zstępujących:
  jeżeli $A_n\in\mathcal{M}$ i $A_1\supseteq A_2\supseteq\cdots$, to
  \[
    \bigcap_{n=1}^{\infty}A_n\in\mathcal{M}.
  \]
\end{enumerate}
\end{definicja}

\begin{przyklad}[Przedziały i zbiory skończone]
Rodzina $\mathcal{J}$ wszystkich przedziałów w $\R$ wraz ze zbiorem
pustym jest rodziną monotoniczną. Przedziały mogą być nieograniczone.
Istotnie, suma wstępującego ciągu przedziałów jest przedziałem
lub zbiorem pustym, podobnie jak przekrój zstępującego ciągu przedziałów.

Natomiast rodzina zbiorów skończonych nie jest monotoniczna.
Połóżmy
\[
  A_n=\{1,2,\ldots,n\}.
\]
Każdy $A_n$ jest skończony i $A_n\subseteq A_{n+1}$, ale
\[
  \bigcup_{n=1}^{\infty}A_n=\N
\]
nie jest zbiorem skończonym.
\end{przyklad}

\begin{uwaga}[Rodziny generowane]
Podobnie jak w przypadku sigma-ciał można rozważać najmniejsze rodziny
danego typu zawierające ustaloną rodzinę zbiorów.
W szczególności przez $\mathcal{M}(\mathcal{E})$ oznaczamy najmniejszą
rodzinę monotoniczną zawierającą
$\mathcal{E}\subseteq\mathcal{P}(X)$.
Jest ona przekrojem wszystkich rodzin monotonicznych zawierających
$\mathcal{E}$:
\begin{equation}
  \mathcal{M}(\mathcal{E})
  =\bigcap\{\mathcal{H}\subseteq\mathcal{P}(X):
    \mathcal{H}\text{ jest monotoniczna, }\mathcal{E}\subseteq\mathcal{H}\}.
  \label{eq:rodzina-monotoniczna-generowana}
\end{equation}
Rodzina pod przekrojem jest niepusta, ponieważ należy do niej
$\mathcal{P}(X)$. Przekrój zachowuje oba warunki
definicji~\ref{def:rodzina-monotoniczna}, co uzasadnia tę konstrukcję.
\end{uwaga}

\begin{twierdzenie}[O rodzinie monotonicznej generowanej przez ciało]
\label{tw:rodzina-monotoniczna}
Jeżeli $\mathcal{F}$ jest ciałem zbiorów w $X$, to
\begin{equation}
  \sigma(\mathcal{F})=\mathcal{M}(\mathcal{F}).
  \label{eq:twierdzenie-monotoniczne}
\end{equation}
\end{twierdzenie}

\begin{proof}
\emph{Pierwsza inkluzja.}
Sigma-ciało $\sigma(\mathcal{F})$ zawiera $\mathcal{F}$ i jest rodziną
monotoniczną, gdyż jest zamknięte na dowolne przeliczalne sumy
i przekroje. Z minimalności $\mathcal{M}(\mathcal{F})$ otrzymujemy
\[
  \mathcal{M}(\mathcal{F})\subseteq\sigma(\mathcal{F}).
\]
Ta inkluzja nie wymaga założenia, że $\mathcal{F}$ jest ciałem.

\emph{Przeciwna inkluzja.}
Wykażemy, że $\mathcal{M}(\mathcal{F})$ jest sigma-ciałem.
Dla skrócenia zapisu połóżmy $\mathcal{M}=\mathcal{M}(\mathcal{F})$.

\emph{1. Zbiór pusty.}
Ponieważ $\mathcal{F}$ jest ciałem, mamy
$\varnothing\in\mathcal{F}\subseteq\mathcal{M}$.

\emph{2. Dopełnienia.}
Rozważmy rodzinę
\[
  \mathcal{G}=\{A\in\mathcal{M}:A^c\in\mathcal{M}\}.
\]
Jest ona podrodziną $\mathcal{M}$ i zawiera $\mathcal{F}$, ponieważ
$\mathcal{F}$ jest zamknięte na dopełnienia.
Pokażemy, że $\mathcal{G}$ jest monotoniczna.
Jeżeli $A_n\in\mathcal{G}$ i $A_1\subseteq A_2\subseteq\cdots$,
to $\bigcup_{n=1}^{\infty}A_n\in\mathcal{M}$ oraz
\[
  \left(\bigcup_{n=1}^{\infty}A_n\right)^c
  =\bigcap_{n=1}^{\infty}A_n^c\in\mathcal{M},
\]
gdyż ciąg $(A_n^c)$ jest zstępujący.
Stąd $\bigcup_{n=1}^{\infty}A_n\in\mathcal{G}$.
Analogicznie dla $A_1\supseteq A_2\supseteq\cdots$ mamy
$\bigcap_{n=1}^{\infty}A_n\in\mathcal{M}$ i
\[
  \left(\bigcap_{n=1}^{\infty}A_n\right)^c
  =\bigcup_{n=1}^{\infty}A_n^c\in\mathcal{M}.
\]
Zatem również $\bigcap_{n=1}^{\infty}A_n\in\mathcal{G}$.
Rodzina $\mathcal{G}$ jest monotoniczna i zawiera $\mathcal{F}$,
więc z minimalności $\mathcal{M}$ wynika
$\mathcal{M}\subseteq\mathcal{G}$. Stąd $\mathcal{G}=\mathcal{M}$,
czyli $\mathcal{M}$ jest zamknięta na dopełnienia.

\emph{3. Sumy dwóch zbiorów.}
Dla ustalonego $B\subseteq X$ połóżmy
\[
  \mathcal{H}_B=\{A\in\mathcal{M}:A\cup B\in\mathcal{M}\}.
\]
Rodzina $\mathcal{H}_B$ jest monotoniczna. Istotnie, dla wstępującego
ciągu $(A_n)$ jej elementów mamy
\[
  \left(\bigcup_{n=1}^{\infty}A_n\right)\cup B
  =\bigcup_{n=1}^{\infty}(A_n\cup B)\in\mathcal{M},
\]
ponieważ ciąg $(A_n\cup B)$ jest wstępujący.
Sama suma $\bigcup_{n=1}^{\infty}A_n$ też należy do $\mathcal{M}$,
więc należy do $\mathcal{H}_B$.
Dla ciągu zstępującego analogicznie korzystamy z równości
\[
  \left(\bigcap_{n=1}^{\infty}A_n\right)\cup B
  =\bigcap_{n=1}^{\infty}(A_n\cup B)\in\mathcal{M}.
\]

Najpierw niech $B\in\mathcal{F}$.
Dla każdego $A\in\mathcal{F}$ mamy
$A\cup B\in\mathcal{F}\subseteq\mathcal{M}$, więc
$\mathcal{F}\subseteq\mathcal{H}_B$.
Z minimalności $\mathcal{M}$ otrzymujemy
$\mathcal{M}\subseteq\mathcal{H}_B\subseteq\mathcal{M}$.
Zatem
\[
  \forall A\in\mathcal{M}\ \forall B\in\mathcal{F}
  \qquad A\cup B\in\mathcal{M}.
\]
Teraz ustalmy dowolny $A\in\mathcal{M}$ i rozważmy
$\mathcal{H}_A$. Ta rodzina jest monotoniczna, a z właśnie
otrzymanego wyniku i przemienności sumy wynika, że zawiera
$\mathcal{F}$. Ponownie minimalność daje
$\mathcal{M}\subseteq\mathcal{H}_A$. Otrzymujemy więc
\[
  \forall A,B\in\mathcal{M}\qquad A\cup B\in\mathcal{M}.
\]
Indukcyjnie $\mathcal{M}$ jest zamknięta na skończone sumy.

\emph{4. Sumy przeliczalne.}
Niech $A_n\in\mathcal{M}$ dla każdego $n\geq1$. Połóżmy
\[
  B_n=\bigcup_{k=1}^{n}A_k.
\]
Wówczas $B_n\in\mathcal{M}$ oraz $B_n\subseteq B_{n+1}$.
Z monotoniczności otrzymujemy
\[
  \bigcup_{n=1}^{\infty}A_n
  =\bigcup_{n=1}^{\infty}B_n\in\mathcal{M}.
\]
Zatem $\mathcal{M}$ jest sigma-ciałem zawierającym $\mathcal{F}$,
więc $\sigma(\mathcal{F})\subseteq\mathcal{M}$.
Obie inkluzje dają równość~\eqref{eq:twierdzenie-monotoniczne}.
\end{proof}
